\section{The rank-one comparison}\label{sec:rankone}

The field $M$ was chosen so that its rank-one values control the nonlinear
term in the action. We now prove that property. The parameters $K,\sigma$
and $\delta$ remain fixed as in Section~\ref{sec:field}.

\begin{proposition}\label{prop:rankone}
For every $B\in\Sym_N$ and every $a\in\R^N$,
\begin{equation}\label{eq:rankone}
 M(B)=aa^{\transpose}\quad\Longrightarrow\quad
 a^{\transpose}(K^2-B^2)a\ge |a|^{2q}.
\end{equation}
\end{proposition}

The proof compares $X_s(B)$ with a matrix that is a function of $B$ alone.
Its tangent map, integrated in $s$, is an average of unitary conjugations.
Trace monotonicity then removes this average from the needed inequality.

\subsection{A tangent inequality on positive matrices}

Divided differences describe the derivative of matrix functional calculus,
as in the work of Daleckii and Krein \cite{DaleckiiKrein1965}; see
\cite[Theorem~2.3.1]{Hiai2010} for a modern finite-dimensional formulation.
We prove the required derivative formula and tangent inequality directly
from the resolvent representation below.

\begin{lemma}\label{lem:resolvent-tangent}
For $y\ge0$ the function $f$ of \eqref{eq:f} has the representation
\begin{equation}\label{eq:resolvent-f}
 f(y)=c_\beta\int_0^\infty v^\beta
 \left(\frac1{\delta+v}-\frac1{y+\delta+v}\right)\,dv,
 \qquad
 c_\beta^{-1}=\int_0^\infty\frac{t^\beta}{(1+t)^2}\,dt.
\end{equation}
If $X,Y\ge0$ are Hermitian matrices, then
\begin{equation}\label{eq:tangent}
 f(X)\le f(Y)+T_Y(X-Y),
\end{equation}
where, in an eigenbasis of $Y$ with eigenvalues $y_i$,
\begin{equation}\label{eq:divdiff}
 (T_Y Z)_{ij}=f[y_i,y_j]Z_{ij},\qquad
 f[x,y]=\int_0^1 f'(\tau x+(1-\tau)y)\,d\tau.
\end{equation}
For $x\ne y$, this is $(f(x)-f(y))/(x-y)$; for $x=y$ it is $f'(x)$.
\end{lemma}
\begin{proof}
The integrand in \eqref{eq:resolvent-f} is $O(v^\beta)$ near zero and
$O(v^{\beta-2})$ at infinity. The range $-1<\beta<1$ suffices for convergence.
Differentiation gives
\[
 c_\beta\int_0^\infty\frac{v^\beta}{(y+\delta+v)^2}\,dv
 =(y+\delta)^{\beta-1}=f'(y),
\]
and both sides of \eqref{eq:resolvent-f} vanish at zero.

Put $A=Y+(\delta+v)I$ and $Z=X-Y$. Since $A$ and $A+Z$ are positive
definite, the resolvent identity gives
\[
 -(A+Z)^{-1}
 =-A^{-1}+A^{-1}ZA^{-1}
   -A^{-1}Z(A+Z)^{-1}ZA^{-1}
 \le -A^{-1}+A^{-1}ZA^{-1}.
\]
Integrating proves \eqref{eq:tangent}. The coefficient multiplying $Z_{ij}$
in its linear term is
\[
 c_\beta\int_0^\infty
 \frac{v^\beta}{(y_i+\delta+v)(y_j+\delta+v)}\,dv=f[y_i,y_j].
\]
This also covers repeated eigenvalues. All these integrals converge
absolutely, since $\delta>0$.
\end{proof}

\subsection{The integrated tangent map}

Suppose that $M(B)=aa^{\transpose}$ and $m=|a|^2>0$. By
Lemma~\ref{lem:jensen}, for every unit $e$,
\[
 0\le e^{\transpose}M(B)e
 \le\mu\bigl(e^{\transpose}K^2e-(e^{\transpose}Be)^2\bigr).
\]
Strict increase of $\mu$ implies
$e^{\transpose}K^2e\ge(e^{\transpose}Be)^2$. Consequently
\begin{equation}\label{eq:X-positive}
 e^{\transpose}X_s(B)e
 =(s-e^{\transpose}Be)^2
   +e^{\transpose}K^2e-(e^{\transpose}Be)^2\ge0
\end{equation}
for every real $s$: the matrices $X_s(B)$ are positive semidefinite.

To apply the tangent inequality, we compare $X_s(B)$ with
$(sI-B)^2+dI$, which is diagonal in an eigenbasis of $B$.
Choose $d>0$ so that its integrated scalar value is $\mu(d)=m$:
\begin{equation}\label{eq:D-d-Y}
 D=K^2-B^2,\qquad
 d=(m+\delta^\sigma)^{1/\sigma}-\delta,\qquad
 Y_s=(sI-B)^2+dI.
\end{equation}
This choice matches the nonzero eigenvalue of $M(B)$.
When we pair the integrated tangent inequality with $M(B)$,
the identity $M(B)^2=mM(B)$ will cancel the contribution $mI$
against the left side.
We have $X_s(B)=Y_s+(D-dI)$; we do not assert that $D\ge0$.

Apply Lemma~\ref{lem:resolvent-tangent} to $X_s(B)\ge0$ and $Y_s>0$,
and denote the tangent map at $Y_s$ by $T_s$.
Subtract $f(s^2)I$ and integrate symmetrically.
In an eigenbasis of $B$, the shift identity \eqref{eq:shift} gives
$c\int_\R(f(Y_s)-f(s^2)I)\,ds=mI$, with the integral understood
at symmetric cutoffs. Thus
\begin{equation}\label{eq:M-tangent}
 M(B)\le mI+\mu'(d)\mathcal P(D-dI),
 \qquad \mathcal P=\frac{c}{\mu'(d)}\int_\R T_s\,ds.
\end{equation}
The integral of the tangent maps converges absolutely: for large $|s|$,
their matrix coefficients are $O(|s|^{2\beta-2})$.

\begin{lemma}\label{lem:average}
In an eigenbasis of $B$ with eigenvalues $z_i$, the map $\mathcal P$ is
\begin{equation}\label{eq:P-kernel}
 (\mathcal P Z)_{ij}=P_{ij}Z_{ij},\qquad
 P_{ij}=\int_0^1
 \left(1+\frac{\tau(1-\tau)(z_i-z_j)^2}{d+\delta}\right)^{\sigma-1}\,d\tau.
\end{equation}
There is an even probability measure $\nu$ on $\R$ such that
\begin{equation}\label{eq:P-average}
 \mathcal P(Z)=\int_\R e^{itB}Ze^{-itB}\,d\nu(t).
\end{equation}
In particular $\mathcal P(I)=I$, and $\mathcal P$ is self-adjoint for the
trace pairing on Hermitian matrices.
\end{lemma}
\begin{proof}
The eigenvalues of $Y_s$ are $y_i=(s-z_i)^2+d$. Formula
\eqref{eq:divdiff} yields
\[
 f[y_i,y_j]=\int_0^1
 \bigl((s-\tau z_i-(1-\tau)z_j)^2+d+\delta
       +\tau(1-\tau)(z_i-z_j)^2\bigr)^{\beta-1}\,d\tau.
\]
These scalar integrands are nonnegative. Tonelli's theorem and a shift in
$s$ evaluate their integral by the same calculation as $\mu'(d)$ in
Lemma~\ref{lem:mu}, giving \eqref{eq:P-kernel}.

For $h\ge0$ and $x\in\R$, the gamma integral gives
\begin{equation}\label{eq:gamma-mixture}
 (1+hx^2)^{\sigma-1}
 =\frac1{\Gamma(1-\sigma)}
   \int_0^\infty v^{-\sigma}e^{-v}e^{-vhx^2}\,dv.
\end{equation}
The density $v^{-\sigma}e^{-v}/\Gamma(1-\sigma)$ has total mass one.
For fixed $v,h$, $e^{-vhx^2}$ is the characteristic function of a centered
Gaussian with variance $2vh$; when $h=0$, use the point mass at zero.
Choose $\tau$ uniformly in $[0,1]$ and put $h=\tau(1-\tau)/(d+\delta)$.
The resulting mixture of these symmetric Gaussian measures is an even
probability measure $\nu$, independent of $i,j$. Its characteristic
function at $z_i-z_j$ is precisely $P_{ij}$, which proves
\eqref{eq:P-average}. The diagonal coefficients are one. Finally, changing
$t$ to $-t$ in \eqref{eq:P-average} and using cyclicity of trace proves
self-adjointness.
\end{proof}

The integrated tangent inequality \eqref{eq:M-tangent} controls
$\mathcal P(D)$, while the desired conclusion involves $D$ itself. The
next step compares them in precisely the trace pairing with $M(B)$.

\subsection{Trace monotonicity and contraction}

\begin{lemma}\label{lem:trace-average}
With $B,D$ and $\mathcal P$ as above,
\begin{equation}\label{eq:trace-average}
 \tr(M(B)D)\ge\tr(M(B)\mathcal P(D)).
\end{equation}
\end{lemma}
\begin{proof}
We use an elementary instance of the generalized Klein trace inequality
\cite[Proposition~3, p.~289]{Petz1994}: for any increasing real function $g$ and
Hermitian matrices $A_1,A_2$,
\begin{equation}\label{eq:trace-monotonicity}
 \tr\bigl((g(A_1)-g(A_2))(A_1-A_2)\bigr)\ge0.
\end{equation}
Indeed, in eigenbases of the two matrices the trace equals
\[
 \sum_{i,j}(g(\lambda_i)-g(\nu_j))(\lambda_i-\nu_j)
       |\langle v_i,w_j\rangle|^2\ge0.
\]
Apply this to $g=f$, $A_1=UX_s(B)U^*$ and $A_2=X_s(B)$, with
$U=e^{itB}$. Since $U$ commutes with $(sI-B)^2$, expansion of the trace gives
\begin{equation}\label{eq:trace-fixed-s}
 0\le\tr\bigl(f(X_s(B))(2D-UDU^*-U^*DU)\bigr).
\end{equation}
The coefficient $2D-UDU^*-U^*DU$ has trace zero. We may therefore subtract
$f(s^2)I$ inside the trace, integrate over $[-R,R]$, and pass to the
symmetric-cutoff limit defining $M$. The result is
\[
 0\le\tr\bigl(M(B)(2D-UDU^*-U^*DU)\bigr).
\]
This passage is also uniform in $t$: the coefficient has norm at most
$4\|D\|_{\HS}$, whereas $M_R(B)\to M(B)$ in matrix norm.
Average with respect to the even measure $\nu$ from Lemma~\ref{lem:average}.
Both conjugation terms average to $\mathcal P(D)$, proving the claim.
\end{proof}

\begin{proof}[Proof of Proposition~\ref{prop:rankone}]
If $a=0$, both sides of \eqref{eq:rankone} vanish. Otherwise the preceding
construction applies. Multiply the matrix inequality
\eqref{eq:M-tangent} by the positive semidefinite matrix $M=aa^{\transpose}$
and take its trace. Since $M^2=mM$ and $\mathcal P(I)=I$, we obtain
\[
 m\tr M\le m\tr M+\mu'(d)
      \bigl(\tr(M\mathcal P(D))-d\tr M\bigr).
\]
Here $\mu'(d)>0$. Lemma~\ref{lem:trace-average} therefore gives
\[
 a^{\transpose}Da=\tr(MD)
 \ge\tr(M\mathcal P(D))\ge d\tr M=dm.
\]
For nonnegative $x,y$ and $0<\sigma<1$, $(x+y)^\sigma\le x^\sigma+y^\sigma$.
Apply this with $x=d$, $y=\delta$ to get
$m=(d+\delta)^\sigma-\delta^\sigma\le d^\sigma$.
Thus $dm\ge m^{1+1/\sigma}=|a|^{2q}$, completing the proof.
\end{proof}
